\documentclass[11pt]{amsart}

\usepackage[a4paper,margin=2.3cm]{geometry}
\usepackage[T1]{fontenc}
\usepackage{lmodern}
\usepackage{amsmath}
\usepackage{amssymb}
\usepackage{amsthm}
\usepackage{aliascnt}
\usepackage{mathtools}
\usepackage{microtype}
\microtypesetup{expansion=false}
\usepackage{booktabs}
\usepackage{array}
\usepackage{float}
\usepackage{enumitem}
\usepackage{xcolor}
\usepackage{graphicx}
\usepackage{xurl}
\usepackage{hyperref}
\usepackage[nameinlink,capitalise,noabbrev]{cleveref}

\begin{filecontents*}[overwrite]{induced-saturation-references.bib}
@article{MartinSmith2012,
  author  = {Ryan R. Martin and Jason J. Smith},
  title   = {Induced saturation number},
  journal = {Discrete Mathematics},
  volume  = {312},
  number  = {21},
  pages   = {3096--3106},
  year    = {2012},
  doi     = {10.1016/j.disc.2012.06.015},
  url     = {https://doi.org/10.1016/j.disc.2012.06.015},
  note    = {\url{https://doi.org/10.1016/j.disc.2012.06.015}}
}

@article{BehrensEtAl2016,
  author  = {Sarah Behrens and Catherine Erbes and Michael Santana and
             Derrek Yager and Elyse Yeager},
  title   = {Graphs with induced-saturation number zero},
  journal = {The Electronic Journal of Combinatorics},
  volume  = {23},
  number  = {1},
  pages   = {Paper No. P1.54},
  year    = {2016},
  doi     = {10.37236/5095},
  url     = {https://doi.org/10.37236/5095},
  note    = {\url{https://doi.org/10.37236/5095}}
}

@article{AxenovichCsikos2019,
  author  = {Maria Axenovich and M{\'o}nika Csik{\'o}s},
  title   = {Induced saturation of graphs},
  journal = {Discrete Mathematics},
  volume  = {342},
  number  = {4},
  pages   = {1195--1212},
  year    = {2019},
  doi     = {10.1016/j.disc.2019.01.006},
  url     = {https://doi.org/10.1016/j.disc.2019.01.006},
  note    = {\url{https://doi.org/10.1016/j.disc.2019.01.006}}
}

@misc{FanEtAl2025,
  author        = {Xinyue Fan and Sahab Hajebi and Sepehr Hajebi and
                   Sophie Spirkl},
  title         = {Halfway to induced saturation for even cycles},
  year          = {2025},
  eprint        = {2505.24100},
  archiveprefix = {arXiv},
  primaryclass  = {math.CO},
  doi           = {10.48550/arXiv.2505.24100},
  url           = {https://arxiv.org/abs/2505.24100},
  note          = {arXiv:2505.24100v2 [math.CO],
                   \url{https://arxiv.org/abs/2505.24100}}
}

@misc{FanEtAl2026,
  author        = {Xinyue Fan and Sahab Hajebi and Sepehr Hajebi and
                   Sophie Spirkl},
  title         = {Asymmetric induced saturation},
  year          = {2026},
  eprint        = {2606.24763},
  archiveprefix = {arXiv},
  primaryclass  = {math.CO},
  doi           = {10.48550/arXiv.2606.24763},
  url           = {https://arxiv.org/abs/2606.24763},
  note          = {arXiv:2606.24763v1 [math.CO],
                   \url{https://arxiv.org/abs/2606.24763}}
}

@article{ConderDobcsanyi2002,
  author  = {Marston Conder and Peter Dobcs{\'a}nyi},
  title   = {Trivalent symmetric graphs on up to 768 vertices},
  journal = {Journal of Combinatorial Mathematics and Combinatorial Computing},
  volume  = {40},
  pages   = {41--63},
  year    = {2002},
  url     = {https://combinatorialpress.com/jcmcc-articles/volume-040/trivalent-symmetric-graphs-on-up-to-768-vertices/},
  note    = {\url{https://combinatorialpress.com/jcmcc-articles/volume-040/trivalent-symmetric-graphs-on-up-to-768-vertices/}}
}

@article{ConderPotocnik2026,
  author  = {Marston D. E. Conder and Primo{\v{z}} Poto{\v{c}}nik},
  title   = {Edge-transitive cubic graphs: analysis, cataloguing and enumeration},
  journal = {Journal of Algebra},
  volume  = {685},
  pages   = {703--737},
  year    = {2026},
  doi     = {10.1016/j.jalgebra.2025.07.035},
  url     = {https://doi.org/10.1016/j.jalgebra.2025.07.035},
  note    = {\url{https://doi.org/10.1016/j.jalgebra.2025.07.035}}
}

@misc{GraphSymCAT1823,
  author       = {Marston D. E. Conder and Primo{\v{z}} Poto{\v{c}}nik},
  title        = {{CAT(182,3)} description},
  howpublished = {Foster Census, GraphSym},
  url          = {https://fostercensus.graphsym.net/files/CAT-descriptions/CAT_182_3.html},
  note         = {Accessed 25 July 2026,
                  \url{https://fostercensus.graphsym.net/files/CAT-descriptions/CAT_182_3.html}}
}

@article{Coxeter1950,
  author  = {H. S. M. Coxeter},
  title   = {Self-dual configurations and regular graphs},
  journal = {Bulletin of the American Mathematical Society},
  volume  = {56},
  number  = {5},
  pages   = {413--455},
  year    = {1950},
  doi     = {10.1090/S0002-9904-1950-09407-5},
  url     = {https://doi.org/10.1090/S0002-9904-1950-09407-5},
  note    = {\url{https://doi.org/10.1090/S0002-9904-1950-09407-5}}
}

@article{HarutyunyanEtAl2016,
  author  = {Ararat Harutyunyan and Reza Naserasr and Mirko Petru{\v{s}}evski
             and Riste {\v{S}}krekovski and Qiang Sun},
  title   = {Mapping planar graphs into the {Coxeter} graph},
  journal = {Discrete Mathematics},
  volume  = {339},
  number  = {2},
  pages   = {839--849},
  year    = {2016},
  doi     = {10.1016/j.disc.2015.10.007},
  url     = {https://doi.org/10.1016/j.disc.2015.10.007},
  note    = {\url{https://doi.org/10.1016/j.disc.2015.10.007}}
}
\end{filecontents*}

\hypersetup{
  pdftitle={Two Computationally Certified Induced-Saturated Graphs for C12 and C14},
  pdfauthor={Romeo Mugnier de Almeida},
  colorlinks=true,
  linkcolor=blue!45!black,
  citecolor=green!35!black,
  urlcolor=blue!55!black
}
\urlstyle{same}

\newcommand{\Ck}[1]{C_{#1}}
\newcommand{\Pk}[1]{P_{#1}}
\newcommand{\Z}{\mathbb{Z}}
\newcommand{\Gtwelve}{G_{12}}
\newcommand{\Gfourteen}{G_{14}}
\newcolumntype{L}[1]{>{\raggedright\arraybackslash}p{#1}}

\newtheorem{theorem}{Theorem}[section]
\newaliascnt{lemma}{theorem}
\newtheorem{lemma}[lemma]{Lemma}
\aliascntresetthe{lemma}
\crefname{lemma}{Lemma}{Lemmas}
\Crefname{lemma}{Lemma}{Lemmas}
\newaliascnt{proposition}{theorem}
\newtheorem{proposition}[proposition]{Proposition}
\aliascntresetthe{proposition}
\crefname{proposition}{Proposition}{Propositions}
\Crefname{proposition}{Proposition}{Propositions}
\newaliascnt{corollary}{theorem}
\newtheorem{corollary}[corollary]{Corollary}
\aliascntresetthe{corollary}
\crefname{corollary}{Corollary}{Corollaries}
\Crefname{corollary}{Corollary}{Corollaries}
\newtheorem{maintheorem}{Theorem}
\renewcommand{\themaintheorem}{\Alph{maintheorem}}
\theoremstyle{definition}
\newaliascnt{definition}{theorem}
\newtheorem{definition}[definition]{Definition}
\aliascntresetthe{definition}
\crefname{definition}{Definition}{Definitions}
\Crefname{definition}{Definition}{Definitions}
\theoremstyle{remark}
\newaliascnt{remark}{theorem}
\newtheorem{remark}[remark]{Remark}
\aliascntresetthe{remark}
\crefname{remark}{Remark}{Remarks}
\Crefname{remark}{Remark}{Remarks}
\crefname{maintheorem}{theorem}{theorems}
\Crefname{maintheorem}{Theorem}{Theorems}

\setlist[enumerate]{leftmargin=2.35em}
\setlist[itemize]{leftmargin=2.1em}
\emergencystretch=2em
\raggedbottom

\title[Induced-saturated graphs for \(\Ck{12}\) and \(\Ck{14}\)]
{Two Computationally Certified Induced-Saturated Graphs\\
for \(\Ck{12}\) and \(\Ck{14}\)}
\author{Romeo Mugnier de Almeida}
\date{Preliminary version, 25 July 2026}
\keywords{induced saturation, even cycle, Coxeter graph, Foster Census,
computational certificate}

\begin{document}

\begin{abstract}
We give computer-assisted proofs for two explicit induced-saturated graphs.
The graph specified by a released \(182\)-vertex edge list and isomorphic to
the Foster Census graph \(\mathrm{CAT}(182,3)\) is
\(\Ck{12}\)-induced-saturated, and the Coxeter graph in its Fano-plane model is
\(\Ck{14}\)-induced-saturated.  An elementary certificate lemma reduces both
claims to finite statements about cycles with unique chords and induced paths
with prescribed endpoints.  Machine-readable certificates give one witness
for every edge and nonedge, while separately written exhaustive programs
recompute the relevant cycle and path data.  For the Coxeter graph, a further
C++ program examines all \(\binom{28}{14}=40{,}116{,}600\) vertex subsets to
check the negative condition.  These computations are not proof-assistant
formalizations.  We make no claim of priority or minimality.
\end{abstract}

\maketitle

\section{Introduction}\label{sec:introduction}

For a fixed graph \(H\), an \(H\)-induced-saturated graph contains no induced
copy of \(H\), while deleting any present edge and adding any missing edge
each creates an induced copy of \(H\).  Martin and Smith introduced an
induced-saturation number in a generalized-graph framework
~\cite{MartinSmith2012}.  Behrens, Erbes, Santana, Yager, and Yeager later
studied the zero case, which yields ordinary \(H\)-induced-saturated graphs
~\cite{BehrensEtAl2016}; further existence results were obtained by Axenovich
and Csik{\'o}s~\cite{AxenovichCsikos2019}.

Even cycles provide a particularly concrete existence problem.  Fan, Hajebi,
Hajebi, and Spirkl construct induced-saturated graphs for the even cycles
through \(\Ck{10}\), and deletion-saturated graphs for every even cycle
~\cite{FanEtAl2025}.  Their subsequent work records that the existence of
induced-saturated graphs for all even cycles remains open
~\cite{FanEtAl2026}.  The present preliminary note records two finite
constructions for \(\Ck{12}\) and \(\Ck{14}\).

\begin{maintheorem}\label{thm:c12}
The graph defined by the released \(182\)-vertex edge list, which is
isomorphic to \(\mathrm{CAT}(182,3)\), is
\(\Ck{12}\)-induced-saturated.
\end{maintheorem}

\begin{maintheorem}\label{thm:c14}
The Coxeter graph is \(\Ck{14}\)-induced-saturated.
\end{maintheorem}

Both proofs use the same reduction.  For an edge deletion we exhibit a cycle
in which the deleted edge is the unique chord.  For a nonedge addition we
exhibit an induced path on the required number of vertices whose endpoints
are the added pair.  The graph-theoretic claims are thereby reduced to
finite, explicitly enumerated statements.

This preliminary note does not establish priority for either construction.
Its purpose is to record explicit, independently checkable computational
certificates.  No claim of minimality is made.
% TODO before public release: insert permanent repository URL

\section{Definitions and the certificate lemma}\label{sec:certificate}

All graphs are finite, simple, and undirected.  For a graph \(G\), write
\(V(G)\) and \(E(G)\) for its vertex and edge sets.  If \(e\in E(G)\), then
\(G-e\) denotes deletion of \(e\).  If \(uv\notin E(G)\) with \(u\ne v\),
then \(G+uv\) denotes addition of the edge \(uv\).  We write \(\Ck{k}\) for
the cycle on \(k\) vertices and \(\Pk{k}\) for the path on \(k\) vertices.

\begin{definition}
Let \(H\) be a graph.  A graph \(G\) is \emph{\(H\)-induced-saturated} if
\begin{enumerate}[label=\textup{(\roman*)}]
  \item \(G\) contains no induced copy of \(H\);
  \item for every \(e\in E(G)\), the graph \(G-e\) contains an induced copy
        of \(H\); and
  \item for every \(uv\notin E(G)\), the graph \(G+uv\) contains an induced
        copy of \(H\).
\end{enumerate}
\end{definition}

A \emph{chord} of a specified cycle is an edge joining two nonconsecutive
vertices of that cycle.

\begin{lemma}[Certificate lemma]\label{lem:certificate}
Let \(G\) be a graph and let \(k\geq 4\).
\begin{enumerate}[label=\textup{(\alph*)}]
  \item If every \(k\)-cycle in \(G\) has at least one chord, then \(G\)
        contains no induced \(\Ck{k}\).
  \item If every edge \(e\in E(G)\) is the unique chord of some
        \(k\)-cycle, then \(G-e\) contains an induced \(\Ck{k}\).
  \item If every nonedge \(uv\) is the endpoint pair of an induced path on
        \(k\) vertices, then \(G+uv\) contains an induced \(\Ck{k}\).
\end{enumerate}
\end{lemma}

\begin{proof}
An induced \(k\)-cycle is precisely a \(k\)-cycle with no chord, which proves
\textup{(a)}.  For \textup{(b)}, let \(Q\) be a \(k\)-cycle whose only chord
is \(e\).  After deleting \(e\), the subgraph induced by \(V(Q)\) consists
exactly of the cycle edges of \(Q\).  For \textup{(c)}, an induced path on
\(k\) vertices has \(k-1\) edges and no further edges among its vertices.
Adding the edge between its endpoints therefore creates an induced
\(k\)-cycle.
\end{proof}

\begin{remark}\label{rem:evidence-layers}
Stored-certificate validation and independent recomputation are distinct.
The former checks that each listed cycle or path is a valid witness and that
the witness keys give exact toggle coverage.  The latter enumerates the
relevant objects directly from the graph, without using the stored witness
chosen for a particular toggle.
\end{remark}

\section{A \texorpdfstring{\(\Ck{12}\)}{C12}-induced-saturated graph}
\label{sec:c12}

\subsection{Construction and graph identity}

Let \(\Gtwelve\) be the released edge-list graph.  Its vertex set is
\(\{0,1,\ldots,181\}\), and its edges are listed in
\texttt{CAT182\_3.edgelist}.  The strict verifier reconstructs this graph and
obtains
\[
  |V(\Gtwelve)|=182,\qquad |E(\Gtwelve)|=273,\qquad
  \binom{182}{2}-273=16{,}198
\]
nonedges.  It also checks that every vertex has degree \(3\) and that the
graph is connected.

Graph identity is checked separately from induced saturation.  The
certificate contains an LCF shift sequence whose reconstructed edge set is
exactly the released edge-list graph.  The direct GraphSym description
accessed on 25 July 2026 lists the same shift sequence for
\(\mathrm{CAT}(182,3)\)~\cite{GraphSymCAT1823}.  An alternative Foster Census
LCF description included in the bundle is textually different; a separate
NetworkX \texttt{GraphMatcher} computation reconstructs both graphs and
returns an isomorphism with \(182\) mapping entries.  Thus the claim made here
is that the edge-list graph is \emph{isomorphic to}
\(\mathrm{CAT}(182,3)\); no equality of differently labelled census
representations is asserted.  For census context and the current catalogue,
see~\cite{ConderDobcsanyi2002,ConderPotocnik2026}.

\subsection{Cycles and edge deletions}

\begin{proposition}\label{prop:c12-cycles}
The graph \(\Gtwelve\) has exactly \(273\) simple \(12\)-cycles.  Every such
cycle has exactly one chord, and every edge of \(\Gtwelve\) is the unique
chord of exactly one such cycle.
\end{proposition}

The primary program exhaustively extends simple \(11\)-edge paths on
\(12\) distinct vertices and closes them when the last vertex is adjacent to
the initial vertex.  It fixes the least cycle vertex as the start and breaks reversal
symmetry by comparing the second and last vertices.  For every retained
cyclic order it subtracts the \(12\) ring edges from the full induced edge
set.  The result is the following complete chord distribution.

\begin{table}[htbp]
\centering
\caption{Chord distribution of the simple \(12\)-cycles in \(\Gtwelve\).}
\label{tab:c12-chords}
\begin{tabular}{@{}rr@{}}
\toprule
Number of chords & Number of \(12\)-cycles \\
\midrule
1 & 273 \\
\bottomrule
\end{tabular}
\end{table}

A separately written standard-library verifier instead generates oriented
simple \(12\)-cycles from every start and canonicalizes only after closure by
rotation and reversal.  It reproduces the \(273\)-cycle count, the chord
histogram, and multiplicity one on every edge.  It also performs a direct
incremental chord-rejecting search in the original graph and finds no induced
\(\Ck{12}\).  For every edge \(e\), a further direct search in
\(\Gtwelve-e\), without reading the stored witness for \(e\), finds an
induced \(\Ck{12}\).

The certificate separately stores one cyclic order for every edge.  Its
validator checks distinct and in-range vertices, all ring edges, the absence
of unintended induced edges, and that the indexed edge is the unique chord.
All \(273\) stored deletion witnesses pass.

\subsection{Induced paths and nonedge additions}

\begin{proposition}\label{prop:c12-paths}
The graph \(\Gtwelve\) has exactly \(218{,}946\) induced paths on
\(12\) vertices, counted up to reversal.  The set of endpoint pairs of these
paths is precisely the set of all \(16{,}198\) nonedges of \(\Gtwelve\).
\end{proposition}

The exhaustive path search starts at every vertex and extends through every
unused neighbour of the current endpoint.  A new endpoint is accepted
exactly when it has no neighbour among the earlier path vertices other than
the current endpoint.  This condition is necessary and sufficient to
preserve inducedness.  At \(12\) vertices, the program retains the orientation
whose initial endpoint has smaller label.  Both implementations obtain
\(218{,}946\) paths and exact coverage of the \(16{,}198\) nonedges.  This
second computation proves every addition case without using the stored
addition witnesses; it is an exhaustive induced-path endpoint computation,
not a separate generic induced-cycle search in each augmented graph.

The certificate stores one induced \(12\)-vertex path for each nonedge.
For every entry the strict validator checks \(12\) distinct vertices, the
\(11\) consecutive path edges, the absence of any other induced edge, and
the exact endpoint pair.  It then checks the toggled induced subgraph.
All \(16{,}198\) stored addition witnesses pass.

\begin{proof}[Proof of \Cref{thm:c12}]
By \cref{prop:c12-cycles}, every \(12\)-cycle in \(\Gtwelve\) has a chord, so
\cref{lem:certificate}\textup{(a)} implies that \(\Gtwelve\) contains no
induced \(\Ck{12}\).  The same proposition states that every edge is the
unique chord of a \(12\)-cycle, so \cref{lem:certificate}\textup{(b)}
implies that every edge deletion creates an induced \(\Ck{12}\).
Finally, \cref{prop:c12-paths} supplies an induced path on \(12\) vertices
between the endpoints of every nonedge.  Hence
\cref{lem:certificate}\textup{(c)} implies that every nonedge addition
creates an induced \(\Ck{12}\).
\end{proof}

\section{The Coxeter graph and \texorpdfstring{\(\Ck{14}\)}{C14}}
\label{sec:c14}

\subsection{Fano-plane construction}

Work in \(\Z_7=\{0,1,\ldots,6\}\), with arithmetic modulo \(7\), and let
\[
  \mathcal L=\bigl\{\{i,i+1,i+3\}:i\in\Z_7\bigr\}.
\]
These are the seven lines of a Fano plane.  Define \(\Gfourteen\) by
\begin{equation}\label{eq:coxeter-fano}
  V(\Gfourteen)=\binom{\Z_7}{3}\setminus\mathcal L,
  \qquad
  XY\in E(\Gfourteen)\quad\Longleftrightarrow\quad X\cap Y=\varnothing .
\end{equation}
This is the standard nonline-triples model of the Coxeter graph
~\cite{Coxeter1950,HarutyunyanEtAl2016}.

There are \(\binom{7}{3}-7=28\) vertices.  We include the cubicity argument.
Fix a nonline triple \(X\).  Each of its three pairs belongs to a distinct
Fano line meeting \(X\) in two points; these three lines account for six of
the nine point-line incidences involving the points of \(X\).  The remaining
three incidences occur on three further lines, each meeting \(X\) in one
point.  Consequently exactly six Fano lines meet \(X\), and exactly one Fano
line is disjoint from \(X\).  The four triples contained in the four-point
complement \(\Z_7\setminus X\) are precisely the possible disjoint triples.
Exactly one is the disjoint Fano line, so the other three are graph
neighbours of \(X\).  Thus \(\Gfourteen\) is cubic, has \(42\) edges, and has
\(\binom{28}{2}-42=336\) nonedges.

The released Python verifiers reconstruct \(\Gfourteen\) from
\cref{eq:coxeter-fano}.  They check exact agreement with the
\(28\)-entry lexicographic triple-label map, the \(42\)-edge file, and the
label map embedded in the certificate.  The primary verifier also recomputes
connectedness, girth \(7\), diameter \(4\), and the unordered-pair distance
distribution
\[
  d_1:42,\qquad d_2:84,\qquad d_3:168,\qquad d_4:84.
\]
The released programs do not compare this graph with an external Coxeter
graph file; the identification as the Coxeter graph comes from the standard
Fano-plane model cited above.

\subsection{Cycles and edge deletions}

\begin{proposition}\label{prop:c14-cycles}
The Coxeter graph has exactly \(420\) simple \(14\)-cycles.  Of these,
\(252\) have one chord and \(168\) have two chords; none has zero chords.
Every edge is the unique chord of exactly six \(14\)-cycles.
\end{proposition}

\begin{table}[htbp]
\centering
\caption{Chord distribution of the simple \(14\)-cycles in the Coxeter graph.}
\label{tab:c14-chords}
\begin{tabular}{@{}rr@{}}
\toprule
Number of chords & Number of \(14\)-cycles \\
\midrule
0 & 0 \\
1 & 252 \\
2 & 168 \\
\bottomrule
\end{tabular}
\end{table}

The primary verifier uses a minimum-rooted depth-first search and breaks
reversal symmetry.  A separately written verifier uses integer adjacency
masks, forms cycles by pairing internally vertex-disjoint \(7\)-edge paths
with common endpoints, and canonicalizes under every rotation and reversal.
Both exhaustive methods give the distribution in
\cref{tab:c14-chords}.  The \(252\) one-chord cycles distribute uniformly:
each of the \(42\) edges is the unique chord six times.

The certificate contains one cyclic order for every edge.  Both Python
programs validate all \(42\) entries as unique-chord witnesses; the primary
program also deletes the indexed chord and tests that the selected induced
subgraph is connected and \(2\)-regular.

\subsection{Induced paths and nonedge additions}

\begin{proposition}\label{prop:c14-paths}
The Coxeter graph has exactly \(5{,}040\) induced paths on \(14\) vertices,
counted up to reversal.  Every nonedge is the endpoint pair of at least one
such path.  The exact endpoint multiplicities are given in
\cref{tab:c14-paths}.
\end{proposition}

\begin{table}[htbp]
\centering
\caption{Endpoint multiplicities for induced \(14\)-vertex paths.}
\label{tab:c14-paths}
\begin{tabular}{@{}rrr@{}}
\toprule
Endpoint distance & Number of nonedges & Paths per pair \\
\midrule
2 & 84  & 4 \\
3 & 168 & 18 \\
4 & 84  & 20 \\
\bottomrule
\end{tabular}
\end{table}

Both complete path searches use the same mathematical extension criterion as
in \cref{prop:c12-paths}, but one stores adjacency sets and the other integer
masks.  They agree on exact nonedge coverage and on
\[
  84\cdot4+168\cdot18+84\cdot20=5{,}040.
\]
The certificate separately stores one path for every nonedge.  Both Python
programs validate the exact endpoints, \(14\) distinct vertices, the
\(13\) path edges, and the absence of any other edge on the selected
vertices.  The primary program also adds the indexed nonedge and checks
connected \(2\)-regularity.  All \(336\) addition witnesses pass.

\begin{proof}[Proof of \Cref{thm:c14}]
By \cref{prop:c14-cycles}, every \(14\)-cycle in the Coxeter graph has a
chord.  Thus \cref{lem:certificate}\textup{(a)} implies that the graph
contains no induced \(\Ck{14}\).  The same proposition states that every edge
is the unique chord of a \(14\)-cycle, so
\cref{lem:certificate}\textup{(b)} implies that every edge deletion creates
an induced \(\Ck{14}\).  Finally, \cref{prop:c14-paths} gives an induced path
on \(14\) vertices between the endpoints of every nonedge.
\Cref{lem:certificate}\textup{(c)} therefore implies that every nonedge
addition creates an induced \(\Ck{14}\).
\end{proof}

\section{Computational verification and reproducibility}
\label{sec:repro}

\subsection{Verification layers}

The released computations keep the following logically different tasks
separate.

For \(\Gtwelve\), the primary program constructs the graph from an embedded
LCF shift sequence and exhaustively enumerates the \(12\)-cycles and induced
\(12\)-vertex paths.  Its essential gates use Python \texttt{assert}, so it
must not be run with \texttt{python -O}.  The separately written verifier
instead reads the edge list and certificate and uses only the Python standard
library.  It rejects duplicate raw JSON object keys and duplicate normalized
toggle keys, checks exact graph and certificate coverage, and uses explicit
runtime checks rather than assertions.  In addition to validating every
stored witness, it implements a different cycle canonicalization, a direct
induced-cycle search in the original graph, direct searches after all
\(273\) edge deletions, and an exhaustive induced-path endpoint computation
covering all \(16{,}198\) additions.  It shares the certificate reduction
with the primary program but imports none of its code.

For \(\Gfourteen\), both Python verifiers begin from
\cref{eq:coxeter-fano}, read all three released data representations, and
validate all \(378\) stored toggle witnesses.  Their JSON readers reject
duplicate raw object keys, noncanonical pair keys, and normalization
collisions.  Neither program contains a correctness-critical
\texttt{assert}; all gates remain active under \texttt{python -O}.  The
second program is a separately written, code-independent implementation
using integer masks, meet-in-the-middle cycle assembly, a different
rotation/reversal canonicalization, and a bit-mask induced-path search.  The
two programs implement the same mathematical reduction and share only the
released data format.

\begin{table}[H]
\centering
\caption{Roles of the computational checks.}
\label{tab:evidence}
\begin{tabular}{@{}L{0.25\textwidth}L{0.67\textwidth}@{}}
\toprule
Check & What it establishes \\
\midrule
Certificate validation &
Every stored order is a valid toggle witness, and the distinct canonical
keys cover exactly all edges or all nonedges. \\
\addlinespace
Independent recomputation &
Cycles, chords, induced paths, and endpoint coverage are recomputed from the
graph without using the stored witness selected for a toggle. \\
\addlinespace
Exhaustive enumeration &
The described finite cycle, path, or fixed-size subset search space is
traversed completely. \\
\addlinespace
Graph identity &
Exact labelled edge-set equality or a separately executed graph-isomorphism
test relates released representations. \\
\addlinespace
File integrity &
SHA-256 digests identify the released bytes; they do not prove any
graph-theoretic statement. \\
\bottomrule
\end{tabular}
\end{table}

\subsection{The all-subsets check}

The C++20 program for \(\Gfourteen\) supplies a third check of the negative
condition.  It reconstructs the graph from the Fano-plane definition and
enumerates exactly
\[
  \binom{28}{14}=40{,}116{,}600
\]
subsets of \(14\) vertices in the original graph.  On each subset it computes
all induced degrees and, if every degree is \(2\), checks connectedness.  A
simple \(14\)-vertex graph is an induced \(\Ck{14}\) exactly when it is
connected and \(2\)-regular.  The program examines the full subset count and
finds none.  This C++ program checks only the absence of an induced
\(\Ck{14}\) in the original graph; it does not check either toggle condition,
the certificate, the edge file, or the label map.

\subsection{Reruns, software, and file integrity}

Fresh runs from clean extractions reproduced the mathematical output of both
core C12 programs and both C14 Python programs.  The C14 C++ source compiled
with C++20 and strict warning flags, and both the optimized and
AddressSanitizer/UndefinedBehaviorSanitizer runs examined all
\(40{,}116{,}600\) subsets without a diagnostic.  The core Python programs
require Python 3; all except the optional C12 graph-identity comparison use
only the standard library.  That optional comparison requires NetworkX.

The optional stored C12 identity log records the SHA-256 digest of an earlier
full census file rather than the digest of the extracted LCF line now
bundled.  A fresh run against the manifest-checked bundled line nevertheless
reproduces the graph-level conclusions: equal order and size, girth \(7\),
diameter \(8\), and an isomorphism with \(182\) mapping entries.  The
isomorphism claim in \cref{sec:c12} relies on this fresh run, not on equality
of the two digest lines.

Both \texttt{SHA256SUMS.txt} manifests pass for every listed payload file.
These hashes detect accidental changes and identify the inputs included in
the released bundles.  They are not evidence that a program was executed and
are not mathematical evidence for induced saturation.  The README files give
the interpreter and compiler commands needed to reproduce the computations.

\section{Concluding remarks}\label{sec:conclusion}

The same short certificate lemma accounts for both constructions: cycles
with a prescribed unique chord control deletions, while induced paths with a
prescribed endpoint pair control additions.  The enumerations reveal further
regularity not needed for the proofs.  In \(\Gtwelve\), unique-chord
\(12\)-cycles correspond one-to-one with edges.  In the Coxeter graph, every
edge occurs six times as a unique chord, and the induced-path multiplicity
depends only on endpoint distance.

The examples do not address minimum order, nor do they by themselves give a
general construction for longer even cycles.  It remains natural to ask
whether comparable certificates occur in other cubic symmetric graphs and
whether automorphism-orbit arguments can compress the released witnesses.
No claim of priority or minimality is made for either graph.

\section*{Acknowledgements}

AI tools assisted in candidate generation and in drafting portions of the
verification code and exposition. The claims recorded here are supported by
the accompanying machine-readable certificates and exhaustive verification
programs described in \cref{sec:repro}. These computations are not formal
proof-assistant developments, and some implement the same mathematical
reduction. This manuscript is a preliminary computational record and has not
been peer reviewed.

\bibliographystyle{amsplain}
\bibliography{induced-saturation-references}

\end{document}